\documentclass[language=english,thesis=bachelor,nosecondtitle,onesided,faculty=cb]{hsmw-thesis}
\type{Internship Report}
\title{Preliminary Investigation of Wash Trading Detection on EtherDelta Using Graph-Based Methods}
\author{Fahim}{Rafique}[B.Sc.]
\submissiondate{2026}[5][17]
\defensedate{2026}
\courseofstudy{Applied Mathematics}
\seminargroup{MA18W1-B}
\examiner[Prof. Dr. rer. nat. Dr.-Ing habil. ]{Florian Zaussinger}
\addexaminer[Dipl. Volkswirt, Dipl. Kfm.]{Mario Oettler}


\newpage
\abstract{Wash trading is a form of market manipulation that artificially inflates trading volume without corresponding changes in asset ownership. In decentralized cryptocurrency markets, detecting such behavior is particularly challenging due to the absence of centralized oversight and the pseudonymous nature of blockchain addresses.

This study applies the WTEYE detection framework to a preprocessed EtherDelta dataset in order to identify and quantify wash trading activity. In addition to the standard methodology, a token-level filtering parameter based on minimum daily occurrence is introduced to examine the sensitivity of wash trading estimates to data selection criteria.

The analysis is conducted by systematically varying the token activity threshold and computing the corresponding wash trading share for each case. The results reveal that wash trading estimates vary significantly across different threshold values. In particular, intermediate activity thresholds produce substantially higher wash trading shares, while both low and high thresholds yield comparatively small values.

These findings indicate that wash trading is not uniformly distributed across tokens, but is instead concentrated within a specific range of activity levels. Furthermore, the results suggest that wash trading measurements are sensitive to data filtering choices, highlighting the importance of careful interpretation of empirical results in studies of decentralized financial markets.}




\begin{document}



\section*{\underline{AI Assistant Usage Disclosure}}

Hochschule Mittweida is open to the use of AI and GenAI and aims to create a space of opportunity that supports and promotes the use of AI systems. The use of generative AI is explicitly desired and will be supported and guided within the scope of available resources.

For examples of correct and incorrect AI Assistant usage, please refer to the original, unabridged version of this form, located at this
\href{https://digital-campus.hs-mittweida.de/2026/02/handreichung-lehren-lernen-und-forschen-mit-generativer-ki/}{link}

\underline{\textit{Use of AI Assistants for Research Purposes}}
\newline
I have used AI Assistant(s) for the purposes of my research as part of this report.
\newline
\newline
Yes  $\boxed{\checkmark}$  \hspace{2cm}   No $\boxed{\phantom{\checkmark}}$
\newline 
\newline  
Explanation:
In this report, AI assistants(ChatGPT 5.5, Gemini 3) are responsibly utilized to support writing and visualization, and code debugging. Specifically by refining written content to meet \\
scientific standards and by effectively troubleshooting implemented algorithms.
\newline
\newline
\newline
I confirm in signing below, that I have reported all usage of AI Assistants for my research, and that the report is truthful and complete.
\newline
\newline

\noindent Chemnitz, 17.05.2026 \hfill Fahim Rafique \\
\rule{\linewidth}{0.5pt}
\noindent Location, Date \hfill Author \\

\clearpage


\chapter{Introduction}
	\section{Cryptocurrency, Blockchain, and Ethereum}		
Cryptocurrencies are digital assets that operate on blockchain systems [1]. A blockchain is a distributed ledger technology that records transactions in a transparent and tamper-resistant manner [1]. Instead of relying on a central authority, transactions are verified and stored across a distributed network of participants. 

This report focuses on the Ethereum blockchain. Ethereum is a decentralized blockchain platform that supports programmable transactions through smart contracts [2]. These smart contracts allow developers to build decentralized applications and financial systems, including decentralized exchanges and ERC20 token systems [2], [17]. ERC20 tokens are blockchain-based digital assets created and transferred on Ethereum. 

Blockchain technology has enabled decentralized financial systems, particularly through cryptocurrencies and decentralized exchanges (DEXs). These systems allow users to trade digital assets directly without relying on centralized intermediaries. While this structure improves transparency and accessibility, it also introduces challenges related to market integrity and market manipulation [6], [7].

	\section{Centralized and Decentralized Exchanges}
Cryptocurrency exchanges can generally be divided into centralized exchanges and decentralized exchanges. Centralized exchanges are operated by companies that manage user accounts, hold customer funds, and internally process transactions [10]. In contrast, decentralized exchanges operate through blockchain-based smart contracts, allowing users to trade assets directly from their wallets without transferring custody to a centralized operator [10], [16].

Decentralized exchanges can also differ in their trading mechanisms. Some use order-book systems, where buyers and sellers submit orders that are matched directly with each other. Others use liquidity pool mechanisms, where trades are executed against pooled assets managed by smart contracts [3].

This report focuses on EtherDelta, an Ethereum-based decentralized exchange that follows an order-book model. EtherDelta is relevant for this report because its trading activity is recorded on the Ethereum blockchain and can therefore be reconstructed and analyzed directly from transaction data. Because transactions on Ethereum are permanently recorded on the blockchain, trading activity can be examined as a network of interactions between addresses [8], [9].

	\section{Wash Trading and Market Manipulation}
A major concern in cryptocurrency markets is market manipulation. One common form is wash trading, where the same trader or a coordinated group of traders repeatedly buys and sells an asset without a meaningful change in ownership. The purpose is to artificially inflate trading volume and create the appearance of market activity [4], [5].
Trading volume is commonly interpreted as a signal of liquidity, demand, and investor interest [6], [18]. Artificially inflating this volume can therefore mislead market participants and distort price formation. Detecting wash trading is important because manipulated trading activity can create false perceptions of market legitimacy and influence investment decisions.
Unlike traditional financial systems, blockchain systems represent participants through pseudonymous addresses rather than identifiable entities. As a result, it is difficult to determine whether multiple addresses belong to the same trader or to coordinated groups. This makes detecting wash trading in decentralized cryptocurrency markets particularly challenging [11], [12].
	\section{Wash Trading Research and Prior Literature}
Concerns regarding fake trading volume in cryptocurrency markets were initially raised through industry reports and investigative articles. Much of this early discussion focused on centralized cryptocurrency exchanges, where reported trading volume may not always reflect genuine market activity. One influential example is the report presented by Bitwise Asset Management to the United States Securities and Exchange Commission, which argued that a substantial portion of reported Bitcoin trading volume was artificially generated [14]. Similar concerns were later discussed in market analysis articles and policy discussions [15].\\
This report focuses specifically on wash trading in decentralized exchanges. Decentralized exchanges are important for this analysis because their trading activity is recorded on public blockchain infrastructure. This makes it possible to reconstruct transaction patterns and study trading behavior at the address level. In this context, “on-chain” refers to transaction data that is publicly recorded on the blockchain itself rather than maintained only by an exchange operator.\\
Academic studies have developed more systematic methods for detecting wash trading in decentralized markets. Victor and Weintraud (2021) studied wash trading on decentralized exchanges such as EtherDelta and IDEX. Their approach models trading activity as graphs and identifies suspicious structures using strongly connected components combined with volume-matching techniques [11].
Cui and Gao (2023) introduced the WTEYE framework, which detects wash trading using ERC20 token transfer data recorded on the Ethereum blockchain. Their method constructs transaction graphs and identifies wash trading through cycle structures and balance-based conditions [12].\\
These approaches identify structural patterns in transaction data rather than relying on identity-based verification. Filtering decisions determine which transactions and tokens are included and therefore influence the results. However, the impact of these choices has not been systematically examined in the existing literature [11], [12].
	\section{Research Focus and Contribution}
This report builds primarily on the WTEYE framework proposed by Cui and Gao (2023). The focus of this report is not to develop a new wash trading detection algorithm, but rather to examine how filtering choices influence the estimated level of wash trading.
Specifically, this report investigates the following research question: \\
\textit{How does varying the minimum daily token occurrence threshold affect the estimated wash trading share in an Ethereum-based decentralized exchange dataset?} \\
To answer this question, the WTEYE framework is applied to a preprocessed EtherDelta dataset. The analysis introduces a token-level filtering parameter and systematically varies the minimum daily token occurrence threshold. For each threshold, the wash trading share is computed and compared. \\
The results suggest that wash trading estimates may be sensitive to filtering choices and may not be uniformly distributed across tokens. In particular, tokens with intermediate activity levels appear to exhibit higher wash trading shares than both low-activity and highly active tokens. \\
The main contribution of this report is therefore a sensitivity analysis of wash trading detection in decentralized cryptocurrency markets. The analysis highlights how methodological assumptions, particularly token-level filtering, can influence empirical measurements of wash trading [11], [12]. \\
\newline
The remainder of this report is structured as follows. Section 2 reviews the background and related literature. Section 3 describes the methodology and experimental setup. Section 4 presents the analysis and results. Section 5 concludes the report and discusses implications and future work.





\chapter{Background}
	\section{Wash Trading in Financial Markets}
Wash trading is a form of market manipulation in which an asset is repeatedly bought and sold without a meaningful change in beneficial ownership. The same trader, or a coordinated group of traders, may simultaneously act as both buyer and seller in order to generate artificial trading activity. The primary objective of wash trading is to inflate reported trading volume and create misleading signals regarding market liquidity, demand, or investor interest [4], [5].\\
In financial markets, trading volume is commonly interpreted as an indicator of market activity and asset popularity. Artificially increasing this volume can therefore distort price discovery and mislead market participants into believing that an asset is more actively traded than it actually is [6], [7]. For this reason, wash trading is widely considered a manipulative trading practice in both traditional and cryptocurrency markets [11], [12].\\
This problem is particularly relevant in cryptocurrency markets, especially on blockchain platforms such as Ethereum, where transactions are recorded under pseudonymous addresses rather than identifiable participants. Ethereum supports decentralized exchanges and token standards such as ERC20, which enable peer-to-peer trading through smart contracts [2], [17]. As a result, it is difficult to determine whether different addresses belong to the same entity or to coordinated groups.\\
Decentralized exchanges (DEXs) further increase this challenge. Although all transactions are publicly observable, there is no central authority to enforce trading rules. Traders can therefore generate artificial activity through repeated self-trades or coordinated exchanges between multiple addresses. For this reason, prior research focuses on identifying patterns in transaction data rather than relying on identity-based verification [11], [12].\\
Decentralized exchanges generally operate using either order-book systems or liquidity pool mechanisms. These two exchange models differ in how trades are executed and how liquidity is provided.\\
In order-book-based exchanges, buyers and sellers submit orders that specify the amount and price at which they want to trade an asset. These orders are stored in an order book until matching counterparties are found. A trade occurs when compatible buy and sell orders are matched. This structure is similar to traditional financial exchanges and allows trading activity to be represented as direct interactions between market participants [10].\\
In liquidity pool-based exchanges, users trade against pools of assets managed by smart contracts rather than directly with other traders [3]. These pools are funded by liquidity providers who deposit tokens into the system. Prices are typically determined automatically through mathematical formulas rather than through direct buyer-seller matching. As a result, trading relationships between participants are less explicit than in order-book systems.\\
This report focuses on EtherDelta, which is an Ethereum-based order-book decentralized exchange [11]. EtherDelta is particularly suitable for graph-based analysis because trading interactions occur directly between addresses and are publicly recorded on the Ethereum blockchain [11].

	\section{Graph-Based Detection Approaches}
To analyze trading behavior, researchers represent transactions as graphs [8], [9]. In this representation, blockchain addresses form nodes and transactions between addresses form directed edges. A directed edge indicates the direction of a transaction, such as a token transfer from one address to another. A sequence of connected transactions forms a path within the graph. These paths allow researchers to study how assets move between groups of addresses over time.\\
Wash trading often appears as cyclic transaction patterns within these graphs [5], [11], [12]. In such cases, assets move through a sequence of addresses and eventually return to the starting point without a meaningful change in ownership. Detecting these cycles is therefore important for identifying potentially manipulative trading behavior [12].\\
Graph-based methods typically follow three steps. First, transaction data is converted into a graph. Second, the graph is simplified by removing noise, such as low-value transactions or isolated nodes. Third, structural patterns—such as cycles or strongly connected components—are analyzed to identify suspicious activity [11].\\
Some graph-based methods analyze strongly connected components (SCCs) to identify cyclic trading structures. An SCC is a group of nodes in which each node can reach every other node through directed paths. In the context of blockchain transactions, this means that assets can move repeatedly among a connected group of addresses through chains of transactions.\\
Victor and Weintraud (2021) use SCCs to detect coordinated trading behavior on decentralized exchanges. The underlying assumption is that wash trading frequently produces circular transaction patterns in which assets repeatedly move among connected addresses without a meaningful transfer of ownership. By identifying SCCs, researchers can isolate clusters of addresses that may participate in coordinated or manipulative trading activity.\\
Different methods emphasize different aspects of this structure. Some focus primarily on connectivity patterns, while others incorporate quantitative constraints, such as balance consistency, to distinguish wash trading from legitimate activity [12].

	\section{Victor and Weintraud (2021)}
Victor and Weintraud (2021) propose a method to detect wash trading on decentralized exchanges using order-book data. They represent trading activity as a directed graph, where nodes correspond to trader addresses and edges represent trades.
The method identifies strongly connected components (SCCs) within this graph. These structures capture circular trading patterns, which are characteristic of wash trading [11].\\
To focus on relevant patterns, the method applies an SCC threshold that filters out small or weakly active components. In their study, only components with sufficient activity are considered. In this context, an occurrence refers to the repeated appearance of trading interactions within the transaction graph. Components with at least 100 occurrences therefore represent trading structures that appear frequently enough to be considered significant for analysis. This threshold is part of the methodology of Victor and Weintraud (2021) and should not be confused with the token-level activity threshold introduced later in this report. The present study applies a different filtering criterion based on the number of daily token occurrences rather than on the activity level of strongly connected components. This filtering step reduces noise and focuses the analysis on clusters that are more likely to represent coordinated trading behavior.
After identifying these components, the method evaluates whether trades satisfy a volume-matching condition. Wash trading assumes that participants generate volume without significantly changing their positions. Therefore, inflows and outflows of assets should approximately balance within a small margin. Victor and Weintraud (2021) use a tolerance of 1\% for this condition.\\
While this approach provides a clear and interpretable framework, it depends on fixed parameter choices. In particular, the use of a fixed SCC threshold may exclude smaller but coordinated groups. The method also does not evaluate how sensitive the results are to these parameters.

	\section{Cui and Gao (2023)}
Cui and Gao (2023) introduce the WTEYE(wash trade eye) framework to detect wash trading using on-chain ERC20 token transfers. Instead of order-book data, this approach uses blockchain transaction records directly.\\
For each token, the method constructs a directed graph where nodes represent addresses and edges represent transfers. Wash trading is identified using both structural and quantitative criteria. Structurally, the method detects cycles in the graph, which indicate repeated trading among a set of addresses. It also considers interactions among neighboring nodes to capture more complex patterns [12].
In addition, WTEYE uses a balance-based condition. The method compares total trading volume with the net change in token balances across participating addresses. Wash trading is characterized by high volume and minimal net balance change. This condition is evaluated using a tolerance parameter, set to 10\% in the study [12].\\
The method also applies preprocessing steps, such as removing small transactions and pruning weakly connected nodes. However, it does not specify exact thresholds for these steps. This makes the filtering process less explicit than in the approach of Victor and Weintraud (2021).\\
Unlike Victor and Weintraud (2021), WTEYE does not rely primarily on explicit SCC filtering. Instead, it focuses on whether trading behavior satisfies balance-based conditions associated with wash trading. This makes the framework more flexible for analyzing token-level transaction patterns across different activity levels.\\
WTEYE provides a flexible framework that captures detailed transaction patterns. However, because its preprocessing steps are not fully parameterized, the effect of filtering choices is not systematically evaluated.

	\section{Comparison and Limitations}
Both Victor and Weintraud (2021) and Cui and Gao (2023) use graph-based methods to detect wash trading, but they differ in several important ways.


%table of comparison between the two algorithms
\begin{table}[h!]
\centering
\small % Reduces font size slightly so the academic text fits comfortably
\begin{tabular}{|| p{4.0cm} | p{5.5cm} | p{6.0cm} ||} 
	\hline                      
	\textbf{Feature} & \textbf{Victor \& Weintraud (2021)} & \textbf{Cui \& Gao (2023) -- WTEYE} \\ [0.5ex] 
	\hline\hline
	Data source & Order-book trading data & On-chain ERC20 transfer data \\ \hline
	Main detection focus & Strongly connected components (SCCs) & Cycles and balance consistency \\ \hline
	Filtering approach & Explicit SCC thresholds & Implicit preprocessing and graph pruning \\ \hline
	Detection condition & Volume matching & Balance-based condition \\ \hline
	Main limitation & Fixed structural thresholds & Effect of filtering choices is not systematically evaluated \\ [1ex]
	\hline
\end{tabular}
\caption{Comparison of Victor \& Weintraud (2021) and Cui \& Gao (2023)}
\label{table:comparison_algorithms}
\end{table}

Table 2.1 shows that the two approaches differ not only in their data sources but also in how suspicious trading behavior is identified. Victor and Weintraud (2021) rely more heavily on explicit graph structures and SCC filtering, whereas WTEYE emphasizes balance-based conditions and flexible transaction analysis. These differences influence which forms of trading behavior are ultimately classified as wash trading.\\
Both methods also assume that wash trading appears as cyclic trading behavior. While this assumption captures many cases, it may not cover all forms of manipulation. In addition, focusing on specific structures or thresholds may introduce bias by excluding relevant behavior.\\
These limitations suggest that wash trading estimates depend not only on detection methods but also on how the data is filtered. Both Victor and Weintraud (2021) and Cui and Gao (2023) analyze large transaction datasets, but neither study systematically evaluates how wash trading estimates vary across tokens with different levels of trading activity.\\
In this report, token activity refers to how frequently a token appears in trading transactions within a given time period, measured through the number of daily token occurrences in the dataset. Tokens that appear infrequently in trading records are considered relatively low-activity tokens, whereas tokens that appear in many transactions are considered highly active tokens.\\
The relationship between token activity and wash trading is not yet well understood. Tokens with very low activity may generate insufficient trading volume to support large-scale manipulation. In contrast, tokens with intermediate levels of activity may be more vulnerable to artificial volume inflation because moderate trading activity can be amplified more easily. Highly active tokens may also exhibit wash trading, but genuine market activity may dominate observed trading volume. Because previous studies do not systematically analyze these differences, it remains unclear how token activity influences wash trading estimates.\\
Because this report focuses on token-level transaction activity and filtering behavior across ERC20 tokens, the WTEYE framework provides a more suitable foundation for the analysis. Its use of on-chain transaction data and token-level transfer graphs allows wash trading estimates to be examined across different levels of token activity. This flexibility makes it more appropriate for analyzing how filtering choices influence wash trading measurements.\\
This report addresses this issue by introducing a token-level activity threshold and systematically analyzing its effect on wash trading estimates.\\


		

	\chapter{Methodology}

This section presents the methodology used to detect and quantify wash trading in decentralized exchange data. The analysis follows the WTEYE framework, which identifies wash trading based on transaction graph structures and balance-based conditions [12]. The method is applied to a preprocessed EtherDelta dataset. In addition, a token-level filtering step is introduced to evaluate how data selection affects the resulting estimates.\\
Several approaches have been proposed to detect wash trading in decentralized markets, including the framework developed by Victor and Weintraud (2021) and the WTEYE method [12]. While both approaches rely on graph-based representations of trading activity, they differ in their treatment of data and filtering assumptions.\\
This study adopts the WTEYE framework for three main reasons. First, WTEYE operates directly on token transfer data, which provides a more detailed representation of trading activity on the blockchain. In contrast, the approach of Victor and Weintraud (2021) relies on order-book data, which is specific to certain exchange designs and less directly aligned with on-chain transaction structures.
Second, WTEYE does not impose strict structural thresholds, such as minimum component size. Instead, it identifies wash trading based on balance-neutral trading behavior, where trading volume is high but net changes in token ownership remain small [12]. This makes the method more flexible and suitable for analyzing tokens with different activity levels.\\
Third, WTEYE supports per-token analysis within defined time windows. This structure aligns directly with the experimental design of this study, where tokens are filtered based on daily activity and analyzed across multiple threshold values.\\
For these reasons, the WTEYE framework provides an appropriate foundation for examining the sensitivity of wash trading estimates to token-level filtering.\\

	\chapter{Analysis}
	\chapter{Discussion}
	\chapter{Conclusion}
	\appendix % Anhang
	

	\chapter{UML-Diagramme}
	...
\end{document}